\documentclass[12pt]{article}
\usepackage[utf8]{inputenc}
\usepackage[T1]{fontenc}
\usepackage[margin=1in]{geometry}
\usepackage{setspace}
\usepackage{xcolor}
\usepackage[colorlinks=true, urlcolor=blue!60!black, linkcolor=blue!60!black]{hyperref}
\setlength{\parindent}{2em}
\setlength{\parskip}{1.25em}
\renewcommand{\baselinestretch}{1.0}
\title{\textit{Medell\'{i}n v. Texas}, 552 U.S. 491 (2008)}
\author{Case Brief}
\date{}
\begin{document}
\maketitle

\section{Case Metadata}

\begin{itemize}
  \item \textbf{Case Name:} \textit{Medell\'{i}n v. Texas}, 552 U.S. 491 (2008) (No. 06-984).
  \item \textbf{Briefing Book Placement:} Chapter 2: Separation of Powers and the Executive (Treaty Power, Self-Execution, and Limits on Unilateral Presidential Lawmaking).
  \item \textbf{Dates:} Argued October 10, 2007; Decided March 25, 2008.
  \item \textbf{Lower Court:} Certiorari to the Texas Court of Criminal Appeals (\textit{Ex parte Medell\'{i}n}, 223 S.W.3d 315 (Tex. Crim. App. 2006)).
  \item \textbf{Vote:} 6--3, affirming the Texas Court of Criminal Appeals.
  \item \textbf{Majority Opinion:} Chief Justice John G. Roberts, Jr., joined by Justices Antonin Scalia, Anthony M. Kennedy, Clarence Thomas, and Samuel A. Alito, Jr.
  \item \textbf{Concurring Opinion:} Justice John Paul Stevens (concurring in the judgment).
  \item \textbf{Dissenting Opinion:} Justice Stephen G. Breyer, joined by Justices David H. Souter and Ruth Bader Ginsburg.
  \item \textbf{Primary Opinion Sources:} \href{https://supreme.justia.com/cases/federal/us/552/491/}{Justia U.S. Supreme Court Center} and \href{https://www.law.cornell.edu/supct/html/06-984.ZS.html}{Cornell Legal Information Institute (LII)}.
\end{itemize}

\section{Facts and Procedural History}

In 1969, the United States ratified the Vienna Convention on Consular Relations (requiring detaining authorities to notify arrested foreign nationals of their right to contact their consulate under Article 36(1)(b)) and its Optional Protocol Concerning the Compulsory Settlement of Disputes (submitting Convention disputes to the compulsory jurisdiction of the International Court of Justice (ICJ)). Under Article 94(1) of the United Nations Charter, each member state ``undertakes to comply with the decision of the [ICJ] in any case to which it is a party'' (\href{https://supreme.justia.com/cases/federal/us/552/491/}{552 U.S. at 499--500}).

Petitioner Jos\'{e} Ernesto Medell\'{i}n, a Mexican national living in the United States since preschool, was arrested on June 29, 1993, for the gang rape and murder of two Houston teenagers. After receiving \textit{Miranda} warnings, Medell\'{i}n signed a written waiver and confessed, but local officers did not advise him of his Vienna Convention right to notify the Mexican consulate. Convicted of capital murder and sentenced to death in Texas state court, Medell\'{i}n first raised his Vienna Convention claim in his initial state postconviction habeas petition. Texas courts---and subsequently a federal district court and the Fifth Circuit---held the claim procedurally defaulted and found no prejudice (\href{https://supreme.justia.com/cases/federal/us/552/491/}{552 U.S. at 500--02}).

While Medell\'{i}n's federal appeal was pending, the ICJ decided \textit{Case Concerning Avena and Other Mexican Nationals (Mex. v. U.S.)}, 2004 I.C.J. 12 (Mar. 31, 2004), holding that the United States had breached Article 36 as to 51 named Mexican nationals, including Medell\'{i}n, and must provide ``review and reconsideration'' of their convictions and sentences regardless of state procedural default rules. After the Supreme Court granted certiorari in \textit{Medell\'{i}n v. Dretke}, 544 U.S. 660 (2005) (\textit{Medell\'{i}n I}), President George W. Bush issued a February 28, 2005 Memorandum to the Attorney General declaring that the United States would ``discharge its international obligations'' under \textit{Avena} ``by having State courts give effect to the decision in accordance with general principles of comity'' (and days later gave notice of U.S. withdrawal from the Optional Protocol). After the Court dismissed \textit{Medell\'{i}n I} as improvidently granted---and held in \textit{Sanchez-Llamas v. Oregon}, 548 U.S. 331 (2006), that the Vienna Convention itself did not displace state default rules---Medell\'{i}n filed a second Texas habeas application invoking \textit{Avena} and the President's Memorandum. The Texas Court of Criminal Appeals dismissed it as an abuse of the writ (223 S.W.3d 315), and the Supreme Court again granted certiorari (\href{https://www.law.cornell.edu/supct/html/06-984.ZS.html}{552 U.S. at 502--04}).

\section{Issue(s) Presented}

\begin{enumerate}
  \item Is the ICJ's judgment in \textit{Avena}, pursuant to the Optional Protocol, the U.N. Charter, and the ICJ Statute, directly enforceable as self-executing domestic federal law that preempts state procedural default rules?
  \item Does President Bush's February 28, 2005 Memorandum independently compel state courts to give domestic legal effect to \textit{Avena} and set aside state limitations on successive habeas petitions?
\end{enumerate}

\section{Holding}

No on both issues. In a 6--3 decision authored by Chief Justice Roberts, the Supreme Court held that (1) while \textit{Avena} creates an international law obligation, neither \textit{Avena} nor Article 94 of the U.N. Charter is self-executing domestic federal law enforceable in U.S. courts absent implementing legislation by Congress; and (2) under Category Three of Justice Jackson's framework in \href{https://supreme.justia.com/cases/federal/us/343/579/}{\textit{Youngstown Sheet \& Tube Co. v. Sawyer}}, 343 U.S. 579 (1952), the President lacks unilateral constitutional authority to convert a non-self-executing treaty into binding domestic law that preempts state procedural default rules (\href{https://supreme.justia.com/cases/federal/us/552/491/}{552 U.S. at 498--99, 525--32}).

\section{Reasoning}

Chief Justice Roberts divided the Court's analysis into treaty self-execution and presidential lawmaking power:

\begin{enumerate}
  \item \textbf{Non-Self-Execution of the \textit{Avena} Judgment:} Relying on Chief Justice Marshall's distinction in \textit{Foster v. Neilson}, 27 U.S. (2 Pet.) 253 (1829), Roberts explained that a treaty creates binding domestic law under the Supremacy Clause only if Congress enacts implementing legislation or the treaty's text conveys an intention that it operate of its own force as ``self-executing.'' Parsing the relevant texts, the Court concluded that (a) the Optional Protocol is merely a bare grant of ICJ jurisdiction saying nothing about domestic enforcement; (b) U.N. Charter Article 94(1)'s phrase ``undertakes to comply'' is a commitment for future action by the political branches, not a directive to domestic courts; (c) Article 94(2) confirms this reading by supplying a solely diplomatic enforcement remedy---referral to the U.N. Security Council, where the United States holds a veto---which automatic domestic enforceability would erase; and (d) the ICJ Statute limits ICJ parties to nation-states, not individuals. The Executive Branch's consistent treaty interpretation and the post-ratification practice of all 47 Optional Protocol signatories---none of which treated ICJ judgments as automatically binding in domestic courts---reinforced this conclusion (\href{https://supreme.justia.com/cases/federal/us/552/491/}{552 U.S. at 504--23}).

  \item \textbf{Rejection of Unilateral Executive Lawmaking (\textit{Youngstown} Category Three):} Although the President's foreign-policy interests in reciprocal consular compliance were ``plainly compelling,'' presidential authority ``must stem either from an act of Congress or from the Constitution itself'' (\href{https://supreme.justia.com/cases/federal/us/343/579/}{\textit{Youngstown}}, 343 U.S. at 585). Because a non-self-executing treaty is ratified on the understanding that it lacks domestic legal force without congressional legislation, the President's unilateral effort to give \textit{Avena} binding domestic effect conflicted with the Senate's implicit understanding and fell into \textit{Youngstown} Category Three---where presidential power is at its ``lowest ebb.'' The responsibility to transform a non-self-executing treaty into domestic law belongs exclusively to Congress under Article I, \S~7 (\href{https://supreme.justia.com/cases/federal/us/552/491/}{552 U.S. at 523--30}).

  \item \textbf{Inapplicability of Claims-Settlement Precedents and the Take Care Clause:} The Court rejected the government's appeal to independent foreign-affairs authority under \href{https://supreme.justia.com/cases/federal/us/453/654/}{\textit{Dames \& Moore v. Regan}}, 453 U.S. 654 (1981), and \textit{American Insurance Ass'n v. Garamendi}, 539 U.S. 396 (2003). Those decisions rested on a narrow, 200-year, congressionally acquiesced practice of using executive agreements to settle civil claims between U.S. citizens and foreign states---not an unprecedented presidential directive ordering state courts to reopen final criminal judgments. Finally, the Take Care Clause (Art. II, \S~3) could not sustain the Memorandum because \textit{Avena} was not domestic ``Law'' for the President to execute (\href{https://supreme.justia.com/cases/federal/us/552/491/}{552 U.S. at 530--32}).
\end{enumerate}

\section{Concurring and Dissenting Opinions}

\begin{itemize}
  \item \textbf{Justice Stevens (Concurring in the Judgment):} Stevens rejected any presumption against self-execution and agreed with the dissent that the Vienna Convention itself is self-executing, calling the case closer than the majority acknowledged (\href{https://supreme.justia.com/cases/federal/us/552/491/}{552 U.S. at 533--37}). Nevertheless, he concluded that Article 94(1)'s phrase ``undertakes to comply'' most naturally contemplates future action by the political branches, and agreed that the President's Memorandum could not unilaterally bind state courts. Pointing to Oklahoma's voluntary compliance with \textit{Avena}, Stevens urged Texas to shoulder its moral and diplomatic duty to protect the Nation's honor by granting review voluntarily.

  \item \textbf{Justice Breyer (Dissenting, joined by Justices Souter and Ginsburg):} Breyer faulted the majority for demanding an explicit textual statement of self-execution that multilateral treaty drafters rarely include because nations have differing domestic incorporation rules (\href{https://supreme.justia.com/cases/federal/us/552/491/}{552 U.S. at 538--67}). Surveying two centuries of Supremacy Clause precedent from \textit{Ware v. Hylton} (1796) onward, Breyer advocated a practical, context-specific, seven-factor inquiry---emphasizing that the underlying Vienna Convention right is self-executing, the Optional Protocol provides for ``compulsory settlement,'' reviewing criminal-procedure prejudice is quintessential judicial work, and the President supported enforcement without congressional objection. Warning that the majority's rule jeopardized roughly 70 U.S. treaties containing ICJ dispute-settlement clauses, Breyer declined to decide the Part III executive-power issue but expressed doubt that the President lacked all power in \textit{Youngstown} Category Two to implement ratified treaty obligations against state law.
\end{itemize}

\section{Key Quotations}

\begin{quote}
``In sum, while treaties `may comprise international commitments . . . they are not domestic law unless Congress has either enacted implementing statutes or the treaty itself conveys an intention that it be ``self-executing'' and is ratified on these terms.' '' --- Chief Justice Roberts, \href{https://supreme.justia.com/cases/federal/us/552/491/}{552 U.S. at 505}.
\end{quote}

\begin{quote}
``The President has an array of political and diplomatic means available to enforce international obligations, but unilaterally converting a non-self-executing treaty into a self-executing one is not among them. The responsibility for transforming an international obligation arising from a non-self-executing treaty into domestic law falls to Congress.'' --- Chief Justice Roberts, \href{https://supreme.justia.com/cases/federal/us/552/491/}{552 U.S. at 525--26}.
\end{quote}

\begin{quote}
``When the President asserts the power to `enforce' a non-self-executing treaty by unilaterally creating domestic law, he acts in conflict with the implicit understanding of the ratifying Senate. His assertion of authority, insofar as it is based on the pertinent non-self-executing treaties, is therefore within Justice Jackson's third category, not the first or even the second.'' --- Chief Justice Roberts, \href{https://supreme.justia.com/cases/federal/us/552/491/}{552 U.S. at 527}.
\end{quote}

\begin{quote}
``The majority reaches a different conclusion because it looks for the wrong thing (explicit textual expression about self-execution) using the wrong standard (clarity) in the wrong place (the treaty language). Hunting for what the text cannot contain, it takes a wrong turn.'' --- Justice Breyer, dissenting, \href{https://supreme.justia.com/cases/federal/us/552/491/}{552 U.S. at 562}.
\end{quote}

\section{Doctrinal Significance and Modern Legacy (Through 2026)}

\textit{Medell\'{i}n v. Texas} reshaped two pillars of U.S. foreign-relations constitutional law: treaty self-execution and the domestic boundaries of Article II executive power. (Following the decision, Congress declined to enact implementing legislation, and Texas executed Medell\'{i}n in August 2008.)

First, \textit{Medell\'{i}n} cemented a rigorous, text-driven threshold for treaty self-execution, raising the bar for litigants seeking to enforce international obligations or tribunal rulings directly in U.S. courts. Through 2025--2026, federal appellate courts continue to apply \textit{Medell\'{i}n}'s textual framework to distinguish non-self-executing commitments from mandatory judicial directives---for example, in international commercial arbitration disputes where circuits (including the Second Circuit in 2025, joining the First, Fourth, Fifth, and Ninth Circuits) have applied \textit{Medell\'{i}n} to hold that Article II(3) of the New York Convention contains a mandatory directive to courts and therefore operates as self-executing federal law preempting state anti-arbitration insurance statutes (\href{https://www.chaffetzlindsey.com/report/arbitration-in-the-courts-july-2025vol-16/}{Chaffetz Lindsey (2025)}).

Second, on separation of powers, \textit{Medell\'{i}n} stands as a landmark modern limit on presidential unilateralism in foreign affairs. Even though the Court later recognized an exclusive presidential power over formal recognition of foreign sovereigns in \href{https://supreme.justia.com/cases/federal/us/576/1/}{\textit{Zivotofsky v. Kerry}}, 576 U.S. 1 (2015), \textit{Medell\'{i}n} sharply cabined \href{https://supreme.justia.com/cases/federal/us/453/654/}{\textit{Dames \& Moore}}'s historical-gloss doctrine and confirmed that foreign-policy urgency cannot empower the President to displace state law or legislate domestically without Congress. Through 2026, \textit{Medell\'{i}n}'s insistence on bicameral legislative authorization before the Executive alters domestic legal rights continues to resonate across major separation-of-powers controversies, echoed in the Court's rejection of novel assertions of unilateral presidential economic policymaking absent clear congressional delegation (e.g., \href{https://supreme.justia.com/cases/federal/us/607/24-1287/}{\textit{Learning Resources, Inc. v. Trump}}, 607 U.S. \_\_\_ (2026), invalidating unilateral presidential tariffs under IEEPA).

\section{Pedagogical Critiques and Exam Hypotheticals}

\subsection*{Key Critiques}

\begin{itemize}
  \item \textbf{Federalism Versus ``One Voice'' in Foreign Affairs:} \textit{Medell\'{i}n} illustrates a sharp tension in constitutional design: while the Supremacy Clause was adopted to prevent individual states from placing the entire Nation in breach of its treaties, Roberts's strict non-self-execution and \textit{Youngstown} Category Three holdings allow a single state's procedural rule to cause an acknowledged international-law breach unless Congress acts.
  \item \textbf{Textual Clarity in Multilateral Treaties:} As Justice Breyer emphasized, foreign treaty negotiators rarely draft clauses addressing domestic U.S. self-execution, raising the critique that \textit{Medell\'{i}n} effectively creates a presumption against self-execution despite disclaiming any need for ``talismanic words.''
\end{itemize}

\subsection*{Exam Hypothetical and ``Opposite Side'' Contrast}

\textbf{Hypothetical (Unenforceable Under \textit{Medell\'{i}n}):} The United States ratifies a multilateral Environmental Disputes Treaty stating that signatory states ``undertake to give effect'' to binding rulings of an International Environmental Tribunal, with noncompliance referred to a diplomatic council. Congress enacts no implementing statute. After the Tribunal rules that a California coastal zoning law violates the treaty, the President issues a Presidential Directive ordering California courts to enjoin enforcement of the zoning law to honor U.S. treaty obligations. \textit{Analysis:} Under \textit{Medell\'{i}n}, both the Tribunal's ruling and the Presidential Directive fail to preempt California law. The ``undertake to give effect'' text and diplomatic enforcement mechanism mark the treaty as non-self-executing, and under \textit{Youngstown} Category Three, the President cannot unilaterally transform a non-self-executing treaty into binding domestic law overriding state police-power regulations.

\textbf{``Opposite Side'' Contrast (Binding Domestic Law Under Supremacy Clause / \textit{Youngstown} Category One):} Contrast two modifications that reverse the outcome. First, if Congress enacts implementing legislation directing federal and state courts to enforce Tribunal judgments (as Congress did for ICSID arbitral awards in 22 U.S.C. \S~1650a), or if the Senate ratifies a treaty whose text explicitly directs that ``domestic courts of each Party shall enforce Tribunal orders as final judgments,'' the obligation becomes supreme federal law binding California courts under Article VI. Second, if the President resolves a longstanding diplomatic crisis by entering an executive agreement settling commercial civil claims of U.S. citizens against a foreign sovereign within a 200-year tradition of congressional acquiescence, the agreement preempts contrary state civil remedies under \href{https://supreme.justia.com/cases/federal/us/453/654/}{\textit{Dames \& Moore v. Regan}} and \textit{Garamendi}.

\end{document}
